\newpage
\titledquestion{The Fringe of the Matter}

\textbf{The Fringe of the Matter}

Consider shrubs, which store data only at their leaves:
\begin{codeblock}
  datatype 'a shrub = Leaf of 'a | Branch of 'a shrub * 'a shrub
\end{codeblock}

The \textbf{fringe} of a shrub is the sequence of its leaf values, read from left
to right. For instance, both of the following trees have fringe
\code{1, 2, 3}, even though they have completely different shapes:

\begin{center}
  \begin{tikzpicture}[
    level distance=10mm,
    level 1/.style={sibling distance=20mm},
    level 2/.style={sibling distance=11mm},
    every node/.style={font=\small},
    dot/.style={circle, fill, inner sep=1.5pt}
  ]
    \node[dot] {}
      child{node[dot] {}
        child{node {\code{1}}}
        child{node {\code{2}}}
      }
      child{node {\code{3}}}
    ;
  \end{tikzpicture}
  \hspace{2cm}
  \begin{tikzpicture}[
    level distance=10mm,
    level 1/.style={sibling distance=20mm},
    level 2/.style={sibling distance=11mm},
    every node/.style={font=\small},
    dot/.style={circle, fill, inner sep=1.5pt}
  ]
    \node[dot] {}
      child{node {\code{1}}}
      child{node[dot] {}
        child{node {\code{2}}}
        child{node {\code{3}}}
      }
    ;
  \end{tikzpicture}
\end{center}

We are curious whether two shrubs have the same fringe. We will choose to
model the fringe of a shrub as a \textbf{stream} of its leaf values, avoiding
constructing the list of all the leaf values up front. We write
\code{S} to refer to the \code{STREAM} library, whose signature is in the appendix.

You are given \code{append : 'a stream -> 'a stream -> 'a stream}, which appends two streams,
as a helper function.

\begin{parts}
  \part[5]\hfill

    Using \code{append}, implement \code{fringe}.

    \spec
      {fringe}
      {\code{'a shrub -> 'a S.stream}}
      {\code{true}}
      {\code{fringe T} evaluates to the stream of the leaf values of \code{T},
      in left-to-right order.}

    \begin{solutionorbox}[14em]
      \begin{codeblock}
        fun fringe (Leaf x) = S.cons (x, S.empty)
          | fringe (Branch (L, R)) = append (fringe L) (fringe R)
      \end{codeblock}

      The stream's suspended tail is what lets us stop early: the leaves of the
      right subtree are not produced until something actually asks for them.
    \end{solutionorbox}

  \newpage

  \part[3]\hfill

    Suppose we had the following bad implementation of \code{append}:
    \begin{codeblock}
      fun appendBad S1 S2 =
        case S.expose S1 of
          S.Empty => S2
        | S.Cons (x, S1') => S.cons (x, appendBad S1' S2)
    \end{codeblock}

    Is \code{appendBad} maximally lazy? If not, what goes wrong, and why does it
    matter here?

    \begin{solutionorbox}[10em]
      No. It calls \code{S.expose S1} \textbf{before} anyone has requested an
      element of the result, forcing a step of \code{S1}'s computation
      immediately --- so merely \textit{building} the appended stream already
      does work. The \code{S.delay} in the given version prevents this: nothing
      is exposed until the resulting stream is itself exposed.

      This matters here because \code{fringe} appends a stream per
      \code{Branch}. With \code{appendBad}, simply constructing
      \code{fringe T} would force its way down the shrub eagerly, and
      \code{sameFringe} (below) could no longer short-circuit --- it would end up
      producing the fringes it was supposed to avoid building.
    \end{solutionorbox}

  \part[8]\hfill

    Implement \code{sameFringe}. You may find it helpful to define a helper function.

    \spec
      {sameFringe}
      {\code{'a shrub * 'a shrub -> bool}}
      {\code{true}}
      {\code{sameFringe (T1, T2)} $\eeq$ \code{true} if \code{T1} and \code{T2}
      have the same fringe, and \code{false} otherwise.}

    \begin{constraint}
      Your solution must stop on the first mismatched leaf.
    \end{constraint}

    \vspace{3.5cm}

    THIS SECTION LEFT BLANK -- WRITE ANSWER ON NEXT PAGE

    \newpage



    \begin{solutionorbox}[30em]
      \begin{codeblock}
        fun sameFringe (T1, T2) =
          let
            fun walk S1 S2 =
              case (S.expose S1, S.expose S2) of
                (S.Empty, S.Empty) => true
              | (S.Cons (x, S1'), S.Cons (y, S2')) =>
                  x = y andalso walk S1' S2'
              | _ => false
          in
            walk (fringe T1) (fringe T2)
          end
      \end{codeblock}

      The \code{andalso} is what short-circuits: on the first mismatched pair of
      leaves, \code{walk} returns \code{false} without exposing any more of
      either stream. The final \code{\_} case handles fringes of different
      lengths (one stream empty while the other is not).
    \end{solutionorbox}

  \part[3]\hfill

    Why can this be more efficient than flattening both shrubs to lists and
    comparing with \code{=}?

    \begin{solutionorbox}[10em]
      The list-based approach must fully traverse \textbf{both} trees and build
      \textbf{both} complete fringe lists before it can compare anything. The
      stream-based approach produces leaves only on demand and compares them one
      at a time, so as soon as the two fringes first differ it stops --- having
      traversed only as far into each shrub as that first mismatch, and never
      allocating either full fringe.
    \end{solutionorbox}
\end{parts}
